\section{Étapes de conception et validations associées}

\subsection{Mise en place du timer pour la mesure périodique}
Première étape: configurer un timer pour déclencher des mesures toutes les 250 ms.

\subsubsection{Calcul de la configuration}
Comme détaillé précédemment, pour obtenir une période de 250 ms avec une fréquence d'horloge de 84 MHz:
\begin{align*}
P &= \frac{(1 + \text{counter\_period}) \times (1 + \text{prescaler})}{F} \\
0,25 &= \frac{(1 + \text{counter\_period}) \times (1 + \text{prescaler})}{84 \times 10^6} \\
(1 + \text{counter\_period}) \times (1 + \text{prescaler}) &= 21~000~000
\end{align*}

Nous avons choisi:
\begin{itemize}
    \item counter\_period = 65535
    \item prescaler = 319
\end{itemize}

\subsubsection{Validation}
La validation a été effectuée en vérifiant la fréquence de déclenchement des mesures avec l'oscilloscope, confirmant une période de 249,66 ms, très proche des 250 ms visés.

\subsection{Développement du module HC-SR04}
\subsubsection{Principe de fonctionnement}
Le principe du capteur HC-SR04:
\begin{enumerate}
    \item Envoi d'une impulsion de 10 µs sur TRIG
    \item Le capteur émet 8 impulsions ultrasoniques à 40 kHz
    \item Le capteur place ECHO à l'état haut jusqu'au retour de l'onde réfléchie
    \item La durée de l'impulsion ECHO est proportionnelle à la distance
\end{enumerate}

\subsubsection{Mise en œuvre et validation}
\begin{itemize}
    \item Configuration d'un signal PWM très court (Pulse = 3) pour le déclenchement
    \item Utilisation des interruptions externes pour mesurer la durée du signal ECHO
    \item Calcul de la distance: $\text{distance} = \frac{\text{durée} \times 340 \text{ m/s}}{2}$
    \item Validation par comparaison avec une mesure de référence
\end{itemize}

\subsection{Développement du module servomoteur}
\subsubsection{Principe de fonctionnement}
Le servomoteur est contrôlé par un signal PWM de période 20 ms (50 Hz):
\begin{itemize}
    \item Largeur d'impulsion de 1 ms: position 0°
    \item Largeur d'impulsion de 2 ms: position 180°
\end{itemize}

\subsubsection{Mise en œuvre et validation}
\begin{itemize}
    \item Configuration du TIM3 pour générer un signal PWM à 50 Hz
    \item Implémentation d'une fonction pour convertir un pourcentage (0-100\%) en largeur d'impulsion
    \item Validation par observation du mouvement du servomoteur
\end{itemize}

\subsection{Implémentation des modes de fonctionnement}
\subsubsection{Mode 1: Contrôle par distance}
Dans ce mode, la position du servomoteur est proportionnelle à la distance mesurée:
\begin{lstlisting}[caption={Contrôle par distance}, label={lst:mode1}]
float distance = ULTRA_SONIC_GetDistance();
uint8_t percentage = (uint8_t)(distance / 21 * 100);
SERVO_set_servo_percentage(percentage);
\end{lstlisting}

\subsubsection{Mode 2: Contrôle par commande série}
Dans ce mode, la position du servomoteur est définie par une commande reçue via la liaison série.

\subsection{Intégration du système complet}
L'intégration finale a consisté à:
\begin{itemize}
    \item Combiner tous les modules développés
    \item Implémenter la machine à états pour gérer les deux modes
    \item Gérer le changement de mode par bouton-poussoir ou commande série
    \item Transmettre périodiquement la distance mesurée
\end{itemize}
